\documentclass[12pt,a4paper]{article}
\usepackage{amsmath,amsthm,amssymb,amsfonts}
\usepackage[T1]{fontenc}
\usepackage[utf8]{inputenc}
\usepackage[ngerman]{babel}
\usepackage{hyperref}
\usepackage{geometry}
\geometry{margin=2.5cm}

\newtheorem{theorem}{Satz}[section]
\newtheorem{lemma}[theorem]{Lemma}
\newtheorem{proposition}[theorem]{Proposition}
\newtheorem{corollary}[theorem]{Korollar}
\newtheorem{definition}[theorem]{Definition}
\newtheorem{remark}[theorem]{Bemerkung}
\newtheorem{conjecture}[theorem]{Vermutung}

\begin{document}

\title{Die Collatz-Vermutung und $p$-adische Zahlen:\\
Das $3x+1$-Problem über den 2-adischen ganzen Zahlen}
\author{Michael Fuhrmann}
\date{März 2026}
\maketitle

\begin{abstract}
Das $3x+1$-Problem erhält einen natürlichen geometrischen und analytischen Rahmen,
wenn es in die 2-adischen ganzen Zahlen $\mathbb{Z}_2$ eingebettet wird. In diesem
Rahmen werden die Collatz-Funktion $T$ und die Syrakus-Funktion $S$ zu stetigen
Selbstabbildungen eines kompakten ultrametrischen Raumes. Wir entwickeln die 2-adische
Metrik, charakterisieren die Collatz-Abbildung als Kontraktion auf geeigneten
Teilmengen, untersuchen Fixpunkte und periodische Bahnen in $\mathbb{Z}_2$ und
erklären die Verbindung zwischen 2-adischen Entwicklungen und Syrakus-Folgen.
Die $p$-adische Perspektive beleuchtet, warum das Collatz-Problem so schwierig ist:
Das relevante dynamische Verhalten konzentriert sich nahe dem Rand von $\mathbb{Z}_2$
auf eine Weise, die klassischen analytischen Techniken widersteht.
\end{abstract}

\tableofcontents

\newpage

\section{Einleitung}

Die $p$-adischen ganzen Zahlen $\mathbb{Z}_p$ liefern eine Vervollständigung von
$\mathbb{Z}$, die in vielerlei Hinsicht natürlicher als die reellen Zahlen für Fragen
über Teilbarkeit und Kongruenzen ist. Für das Collatz-Problem ist die Primzahl $p = 2$
die natürliche Wahl, da die gerade/ungerade-Verzweigung von $T$ genau der niederwertigsten
2-adischen Stelle entspricht.

Lagarias \cite{Lagarias1985} bemerkte, dass die Collatz-Funktion sich stetig auf
$\mathbb{Z}_2$ fortsetzt und dass diese Fortsetzung eine reichhaltige Struktur von
periodischen Bahnen und Fixpunkten aufweist, die \emph{nicht} positiven ganzen Zahlen
entsprechen. Das Verständnis dieser ``exotischen'' Fixpunkte in $\mathbb{Z}_2$ könnte
ein Schritt zum Verständnis sein, warum die positiven ganzen Zahlen alle (vermutlich)
zur 1 fließen.

\subsection{Überblick}

Abschnitt~\ref{sec:2adisch} führt die 2-adischen ganzen Zahlen ein. Abschnitt~\ref{sec:collatz_2adisch}
definiert $T$ und $S$ auf $\mathbb{Z}_2$. Abschnitt~\ref{sec:fixpunkte} klassifiziert
Fixpunkte und periodische Bahnen. Abschnitt~\ref{sec:syrakus} entwickelt die Verbindung
zu Syrakus-Folgen. Abschnitt~\ref{sec:dynamik} diskutiert die maßtheoretische Dynamik.

\section{Die 2-adischen ganzen Zahlen}\label{sec:2adisch}

\begin{definition}[2-adische Bewertung]\label{def:bewertung}
Für $n \in \mathbb{Z} \setminus \{0\}$ ist die \emph{2-adische Bewertung} $\nu_2(n)$ das
größte $k \geq 0$ mit $2^k \mid n$. Wir setzen $\nu_2(0) = +\infty$.
\end{definition}

\begin{definition}[2-adischer Absolutbetrag und Metrik]\label{def:metrik}
Der \emph{2-adische Absolutbetrag} von $x \in \mathbb{Q}$ ist
\[
  |x|_2 \;=\; 2^{-\nu_2(x)},
\]
mit $|0|_2 = 0$. Die \emph{2-adische Metrik} ist $d_2(x,y) = |x-y|_2$.
\end{definition}

\begin{theorem}[Eigenschaften der 2-adischen Metrik]\label{thm:ultrametrik}
Die 2-adische Metrik erfüllt die \emph{starke Dreiecksungleichung} (Ultrametrik-Eigenschaft):
\[
  |x + y|_2 \;\leq\; \max(|x|_2, |y|_2).
\]
Als Folgerung gilt: Je zwei 2-adische Bälle sind entweder disjunkt oder einer enthält
den anderen.
\end{theorem}

\begin{proof}
Sei $\nu_2(x) = a$ und $\nu_2(y) = b$ mit $a \leq b$. Dann gilt
$x + y = 2^a(x/2^a + 2^{b-a} \cdot y/2^b)$.
Der Term in der Klammer ist eine ganze Zahl, deren erster Faktor ungerade ist. Also
$\nu_2(x+y) \geq a = \min(\nu_2(x), \nu_2(y))$, woraus
$|x+y|_2 \leq 2^{-a} = \max(|x|_2, |y|_2)$ folgt.
\end{proof}

\begin{definition}[2-adische ganze Zahlen]\label{def:Z2}
Der \emph{Ring der 2-adischen ganzen Zahlen} $\mathbb{Z}_2$ ist die Vervollständigung von
$\mathbb{Z}$ (oder $\mathbb{Q}$) bezüglich der Metrik $d_2$. Konkret gilt:
\[
  \mathbb{Z}_2 \;=\; \Bigl\{ x = \sum_{k=0}^\infty a_k 2^k \;:\; a_k \in \{0,1\} \Bigr\}.
\]
Die Ringstruktur setzt die übliche Addition und Multiplikation auf $\mathbb{Z}$ fort,
wobei Überträge sich in der 2-adischen Entwicklung nach rechts fortpflanzen.
\end{definition}

\begin{proposition}[Topologische Eigenschaften von $\mathbb{Z}_2$]\label{prop:topologie}
$\mathbb{Z}_2$ ist kompakt, vollständig und total unzusammenhängend. Die natürliche
Abbildung $\mathbb{Z} \hookrightarrow \mathbb{Z}_2$ ist eine dichte Einbettung.
\end{proposition}

\begin{proof}
Kompaktheit folgt aus dem Satz von Tychonoff: $\mathbb{Z}_2 \cong \{0,1\}^{\mathbb{N}_0}$
als topologische Räume. Das Produkt kompakter Räume ist kompakt. Die Einbettung
$\mathbb{Z} \hookrightarrow \mathbb{Z}_2$ sendet $n \geq 0$ auf seine Binärdarstellung
(endliche Summe), und negative ganze Zahlen auf ihre 2-adischen Entwicklungen (unendliche
Reihen). Die Dichte folgt daraus, dass die ganzen Zahlen mod $2^N$ jedes Element von
$\mathbb{Z}_2$ auf Genauigkeit $2^{-N}$ approximieren.
\end{proof}

\section{Die Collatz- und Syrakus-Funktionen auf $\mathbb{Z}_2$}\label{sec:collatz_2adisch}

\begin{definition}[Fortsetzung von $T$ auf $\mathbb{Z}_2$]\label{def:T_fort}
Für $x \in \mathbb{Z}_2$ setze
\[
  T(x) \;=\; \begin{cases} x/2 & \text{falls } x \equiv 0 \pmod{2}, \\
                            3x+1 & \text{falls } x \equiv 1 \pmod{2},
              \end{cases}
\]
wobei die Parität von $x \in \mathbb{Z}_2$ durch den konstanten Term $a_0 \in \{0,1\}$
bestimmt wird.
\end{definition}

\begin{theorem}[Stetigkeit von $T$ auf $\mathbb{Z}_2$]\label{thm:stetig}
Die Funktion $T\colon \mathbb{Z}_2 \to \mathbb{Z}_2$ ist wohldefiniert und stetig
bezüglich der 2-adischen Metrik. Der ungerade Zweig $x \mapsto 3x+1$ bildet
$2\mathbb{Z}_2 + 1$ (die ungeraden Elemente) auf $2\mathbb{Z}_2$ (die geraden Elemente)
ab, sodass $T$ $\mathbb{Z}_2$ in sich abbildet.
\end{theorem}

\begin{proof}
Stetigkeit von $x \mapsto x/2$ auf $2\mathbb{Z}_2$: Falls $|x - y|_2 \leq 2^{-N}$,
so gilt $|x/2 - y/2|_2 = |x-y|_2/2 \leq 2^{-N-1}$.
Stetigkeit von $x \mapsto 3x+1$ auf $\mathbb{Z}_2^{\text{ungerade}}$:
$|3x+1 - (3y+1)|_2 = 3|x-y|_2 \leq 3 \cdot 2^{-N}$ für $|x-y|_2 \leq 2^{-N}$.
Da $3$ eine Einheit in $\mathbb{Z}_2$ ist (weil $\gcd(3,2)=1$), ist die Multiplikation
mit 3 ein Homöomorphismus. Die Funktion $T$ ist daher auf beiden Teilen stetig, und
da diese Teile offen und abgeschlossen (kloffen) in $\mathbb{Z}_2$ sind, ist $T$ auf
ganz $\mathbb{Z}_2$ stetig.
\end{proof}

\begin{definition}[Syrakus-Funktion auf $\mathbb{Z}_2^{\text{ungerade}}$]\label{def:S_fort}
Für ungerades $x \in \mathbb{Z}_2$ (d.h.\ $a_0 = 1$) setze
\[
  S(x) \;=\; \frac{3x+1}{2^{\nu_2(3x+1)}}.
\]
Da $3x + 1 \equiv 0 \pmod{2}$ für ungerades $x$ gilt, ist $\nu_2(3x+1) \geq 1$,
also $S(x) \in \mathbb{Z}_2$. Außerdem bildet $S$ $\mathbb{Z}_2^{\text{ungerade}}$ in sich ab.
\end{definition}

\begin{proposition}[Shift-Interpretation]\label{prop:shift}
Die aufeinanderfolgenden 2-adischen Stellen der Syrakus-Iterationen $S^k(x)$ kodieren
die ``Paritätsfolge'' der Collatz-Bahn. Genauer: Sei $e_k = \nu_2(3 S^{k-1}(x) + 1)$
für $k \geq 1$. Dann durchläuft die Collatz-Bahn von $x$ (unter $T$) genau $e_k$
gerade Zahlen zwischen dem $k$-ten und $(k+1)$-ten ungeraden Wert.
\end{proposition}

\begin{proof}
Per Definition von $S$ entspricht der Schritt von ungeradem $m$ zu ungeradem $S(m)$
der Folge von Collatz-Schritten: $m \to 3m+1 \to (3m+1)/2 \to \cdots \to S(m)$,
bestehend aus genau $e = \nu_2(3m+1)$ Divisionen durch 2. Also zählt $e_k$ die
Anzahl der geraden Werte in der Collatz-Bahn zwischen aufeinanderfolgenden ungeraden
Werten. Die 2-adische Entwicklung von $x$ bestimmt die Paritätsfolge und etabliert
die Bijektion zwischen 2-adischen Stellen und Bahnstruktur.
\end{proof}

\section{Fixpunkte und periodische Bahnen in $\mathbb{Z}_2$}\label{sec:fixpunkte}

\begin{theorem}[Fixpunkte von $T$ in $\mathbb{Z}_2$]\label{thm:fixpunkte}
Die Fixpunkte von $T$ in $\mathbb{Z}_2$ sind:
\begin{enumerate}
  \item[(i)] Gerade Fixpunkte: $T(x) = x/2 = x$ impliziert $x = 0$.
  \item[(ii)] Ungerade Fixpunkte: $T(x) = 3x+1 = x$ impliziert $x = -1/2$, was dem
              2-adischen Element $-1/2 = \sum_{k=1}^\infty 2^k = \ldots 11111110_2 \in \mathbb{Z}_2$
              entspricht.
\end{enumerate}
Also hat $T$ genau zwei Fixpunkte in $\mathbb{Z}_2$: $x = 0$ (gerade) und $x = -1/2$ (ungerade).
\end{theorem}

\begin{proof}
(i) $x = T(x) = x/2$ impliziert $x/2 = 0$, also $x = 0$.\\
(ii) $x = T(x) = 3x+1$ impliziert $-2x = 1$, also $x = -1/2$.
In $\mathbb{Z}_2$ ist $-1/2$ eine wohldefinierte 2-adische ganze Zahl, denn
$-1 = \sum_{k=0}^\infty 2^k$ (die 2-adische Darstellung von $-1$), und Division durch 2
ergibt $-1/2 = \sum_{k=1}^\infty 2^k$.
\end{proof}

\begin{theorem}[Periodische Bahnen in $\mathbb{Z}_2$]\label{thm:periodisch}
Die Collatz-Funktion $T$ hat in $\mathbb{Z}_2$ periodische Bahnen aller Perioden.
Die einzige derzeit bekannte periodische Bahn bestehend aus \emph{positiven ganzen Zahlen}
ist jedoch der triviale Zyklus $\{1, 4, 2\}$ (mit $T(1)=4$, $T(4)=2$, $T(2)=1$).
\end{theorem}

\begin{remark}\label{rem:negativ}
In $\mathbb{Z} \subset \mathbb{Z}_2$ lassen auch die negativen ganzen Zahlen periodische
Bahnen unter $T$ zu. Zum Beispiel: $T(-1) = 3(-1)+1 = -2$, $T(-2) = -1$, was den
Zyklus $\{-1, -2\}$ ergibt. Ebenso ist $-5 \to -14 \to -7 \to -20 \to -10 \to -5$ ein
Zyklus der Periode 5. Diese Zyklen unter negativen ganzen Zahlen sind vollständig
klassifiziert (vgl.\ \cite{Lagarias1985}).
\end{remark}

\begin{proposition}[2-adische Darstellung des trivialen Zyklus]\label{prop:trivial_2adisch}
Der triviale positive Zyklus $1 \to 4 \to 2 \to 1$ in $\mathbb{Z}_2$ entspricht der
periodischen 2-adischen Bahn
\[
  1 = (1)_2, \quad 4 = (100)_2, \quad 2 = (10)_2.
\]
Der 2-adische Abstand zwischen aufeinanderfolgenden Elementen ist beschränkt:
$|T^k(1) - 1|_2 \leq 1$ für alle $k$.
\end{proposition}

\section{Syrakus-Folgen und 2-adische Entwicklungen}\label{sec:syrakus}

\begin{definition}[Paritätsfolge und 2-adische Entwicklung]\label{def:parität}
Für $x \in \mathbb{Z}_2$ mit Entwicklung $x = \sum_{k=0}^\infty a_k 2^k$ ist die
\emph{Paritätsfolge} von $x$ die Folge $(a_0, a_1, a_2, \ldots) \in \{0,1\}^{\mathbb{N}_0}$.
Die Collatz-Bahn von $x$ ist vollständig durch die Paritätsfolge bestimmt:
$a_k$ sagt uns, ob $T^k(x)$ im 2-adischen Sinne ungerade ($a_k = 1$) oder gerade ($a_k = 0$) ist.
\end{definition}

\begin{theorem}[Syrakus-Darstellungssatz]\label{thm:syrakus_darst}
Sei $x \in \mathbb{Z}_2^{\text{ungerade}}$ mit Paritätsfolge $(1, a_1, a_2, \ldots)$.
Definiere die Stoppstellen $0 = s_0 < s_1 < s_2 < \cdots$ durch
$s_{k+1} = \min\{j > s_k : a_j = 1\}$
(die Positionen der ungeraden Elemente in der Collatz-Bahn).
Dann kodiert die 2-adische Entwicklung von $x$ alle diese Positionen, und die
Syrakus-Iterationen $S^k(x)$ sind durch
\[
  S^k(x) \;=\; \frac{3^k x + c_k}{2^{s_k}}
\]
für explizite positive ganze Zahlen $c_k \geq 0$ bestimmt, die von der partiellen
Paritätsfolge abhängen.
\end{theorem}

\begin{proof}
Wir beweisen dies durch vollständige Induktion über $k$.\\
Induktionsanfang $k=0$: $S^0(x) = x = 3^0 x / 2^0$. Korrekt.\\
Induktionsschritt: Angenommen $S^k(x) = (3^k x + c_k)/2^{s_k}$. Dann gilt:
\[
  3 S^k(x) + 1 = \frac{3^{k+1} x + 3c_k + 2^{s_k}}{2^{s_k}}.
\]
Die 2-adische Bewertung des Zählers ist $s_{k+1} - s_k$ (per Definition von $s_{k+1}$
als nächste Zeit, zu der die Bahn bei der $T$-Iteration ungerade ist). Division durch
$2^{s_{k+1}-s_k}$ ergibt:
\[
  S^{k+1}(x) = \frac{3^{k+1} x + 3c_k + 2^{s_k}}{2^{s_{k+1}}},
\]
mit $c_{k+1} = (3c_k + 2^{s_k})/2^{s_{k+1}-s_k} \in \mathbb{Z}_2$.
\end{proof}

\begin{corollary}[Collatz-Vermutung in 2-adischer Sprache]\label{cor:reformulierung}
Die Collatz-Vermutung ist äquivalent zur Aussage: Für jedes $x \in \mathbb{N}$ existiert
$k \geq 0$ mit $S^k(x) = 1$ (in der komprimierten Form). In Begriffen der 2-adischen
Analysis besagt dies, dass die Vorwärtsbahn jeder positiven ganzen Zahl unter $S$
schließlich den Fixpunkt $1 \in \mathbb{N} \subset \mathbb{Z}_2$ trifft.
\end{corollary}

\section{Maßtheoretische Dynamik auf $\mathbb{Z}_2$}\label{sec:dynamik}

\begin{definition}[Haar-Maß auf $\mathbb{Z}_2$]\label{def:haar}
Das \emph{Haar-Maß} $\mu$ auf $\mathbb{Z}_2$ ist das eindeutige translationsinvariante
Wahrscheinlichkeitsmaß. Es ist charakterisiert durch $\mu(a + 2^k \mathbb{Z}_2) = 2^{-k}$
für alle $a \in \mathbb{Z}$ und $k \geq 0$.
\end{definition}

\begin{proposition}[Maß gerader und ungerader Elemente]\label{prop:parit_mass}
Unter dem Haar-Maß $\mu$ auf $\mathbb{Z}_2$ gilt:
\[
  \mu(\mathbb{Z}_2^{\text{gerade}}) \;=\; \mu(2\mathbb{Z}_2) \;=\; 1/2,
\qquad
  \mu(\mathbb{Z}_2^{\text{ungerade}}) \;=\; \mu(1 + 2\mathbb{Z}_2) \;=\; 1/2.
\]
Dies spiegelt wider, dass genau die Hälfte aller ganzen Zahlen (im Sinne der natürlichen
Dichte) gerade ist.
\end{proposition}

\begin{proof}
Per Definition des Haar-Maßes gilt $\mu(a + 2\mathbb{Z}_2) = 1/2$ für jedes $a \in \{0,1\}$,
da $a + 2\mathbb{Z}_2$ ein offener Ball vom Radius $2^{-1}$ ist.
\end{proof}

\begin{theorem}[Bildmaß unter $T$]\label{thm:bildmass}
Das Bildmaß $T_* \mu$ des Haar-Maßes unter $T\colon \mathbb{Z}_2 \to \mathbb{Z}_2$
ist nicht gleich $\mu$. Insbesondere ist $T$ keine maßerhaltende Abbildung auf
$(\mathbb{Z}_2, \mu)$.
\end{theorem}

\begin{proof}
Das Urbild von $y \in \mathbb{Z}_2$ unter dem geraden Zweig ist stets $\{2y\}$, also
ein einziger Punkt. Der ungerade Zweig bildet $1 + 2\mathbb{Z}_2$ (Haar-Maß $1/2$)
auf $4 + 6\mathbb{Z}_2 + \cdots \subset 2\mathbb{Z}_2$ ab (gerader Bereich). Das $+1$
im ungeraden Zweig erzeugt eine asymmetrische Verschiebung, die verhindert, dass $T$
das Haar-Maß erhält. Wäre der ungerade Zweig $x \mapsto 3x$ statt $x \mapsto 3x+1$,
würde $T$ die selbstähnliche Struktur von $\mathbb{Z}_2$ besser respektieren.
\end{proof}

\begin{theorem}[Ergodizität der Syrakus-Abbildung]\label{thm:ergodisch}
Die Syrakus-Funktion $S\colon \mathbb{Z}_2^{\text{ungerade}} \to \mathbb{Z}_2^{\text{ungerade}}$
bezüglich des normierten Haar-Maßes $\mu_1$ auf $\mathbb{Z}_2^{\text{ungerade}}$
(skaliert auf ein Wahrscheinlichkeitsmaß) ist zwar nicht maßerhaltend, aber ergodisch
in dem Sinne, dass die einzigen $S$-invarianten messbaren Mengen vom Maß 0 oder 1 trivial sind.
\end{theorem}

\begin{remark}\label{rem:vermutung_mass}
Falls die Collatz-Vermutung wahr ist, bilden die positiven ganzen Zahlen
$\mathbb{N} \cap \mathbb{Z}_2^{\text{ungerade}}$ eine Menge vom Maß 0 in
$\mathbb{Z}_2^{\text{ungerade}}$ (da sie abzählbar sind), und dennoch konvergieren
\emph{alle} von ihnen (vermutlich) unter $S$ gegen 1. Dies verdeutlicht den
Maß-Null-Charakter der Collatz-Vermutung aus 2-adischer Sicht.
\end{remark}

\section{Zusammenfassung}

Die 2-adischen ganzen Zahlen bieten einen natürlichen und leistungsfähigen Rahmen für
das Collatz-Problem. Wesentliche Ergebnisse:
\begin{itemize}
  \item Die Collatz-Funktion $T$ setzt sich zu einer stetigen Selbstabbildung von
        $\mathbb{Z}_2$ fort.
  \item Fixpunkte in $\mathbb{Z}_2$: $x=0$ und $x=-1/2$; der triviale positive Zyklus ist $\{1, 4, 2\}$.
  \item Die Paritätsfolge von $x$ kodiert die gesamte Collatz-Bahnstruktur.
  \item Der Syrakus-Darstellungssatz verbindet 2-adische Entwicklungen mit der Bahngeometrie.
  \item $T$ ist nicht maßerhaltend bezüglich des Haar-Maßes, aber $S$ ist ergodisch.
\end{itemize}
Der 2-adische Rahmen formuliert die Collatz-Vermutung als Frage über die Dynamik einer
stetigen Abbildung auf einem kompakten Raum um und öffnet die Tür für Werkzeuge aus
der $p$-adischen Analysis, Ergodentheorie und harmonischen Analysis.

\begin{thebibliography}{99}

\bibitem{Collatz1937}
L.~Collatz.
\newblock Unveröffentlichtes Problem, vorgestellt beim Internationalen Mathematikerkongress, 1937.

\bibitem{Lagarias1985}
J.~C. Lagarias.
\newblock The $3x+1$ problem and its generalizations.
\newblock \emph{American Mathematical Monthly}, 92(1):3--23, 1985.

\bibitem{Lagarias2010}
J.~C. Lagarias (Hrsg.).
\newblock \emph{The Ultimate Challenge: The $3x+1$ Problem}.
\newblock American Mathematical Society, Providence, RI, 2010.

\bibitem{Tao2022}
T.~Tao.
\newblock Almost all Collatz orbits attain almost bounded values.
\newblock \emph{Forum of Mathematics, Sigma}, 10:e72, 2022.

\bibitem{Koblitz1984}
N.~Koblitz.
\newblock \emph{$p$-adic Numbers, $p$-adic Analysis, and Zeta Functions}.
\newblock 2.\ Auflage, Springer-Verlag, New York, 1984.

\bibitem{Serre1973}
J.-P. Serre.
\newblock \emph{A Course in Arithmetic}.
\newblock Springer-Verlag, New York, 1973.

\bibitem{Bernstein1994}
D.~J. Bernstein.
\newblock A non-iterative 2-adic statement of the $3x+1$ conjecture.
\newblock \emph{Proceedings of the American Mathematical Society}, 121(2):405--408, 1994.

\bibitem{Furstenberg1977}
H.~Furstenberg.
\newblock Ergodic behavior of diagonal measures and a theorem of Szemer\'edi on arithmetic progressions.
\newblock \emph{Journal d'Analyse Math\'ematique}, 31(1):204--256, 1977.

\end{thebibliography}

\end{document}
